% GENERATED FILE. Edit canonical lessons, then run npm run generate:books.
% canonical-source-hash: fnv1a64:63ff8f05655588b6
% canonical-lessons: HI-C62-rope, HI-C62-needle, HI-C62-basket, HI-C62-clay, HI-C62-wood

\chapter{What Hands Work With}
\label{ch:hi-what-hands-work-with}
\hlchaptermodality{62}

\begin{chapteropening}
\textbf{By the end of this chapter:} \emph{I can name five things a pair of hands works with when I say what something is called.}

Complete HI-C62-wood and say all five words in order.
\end{chapteropening}

\section[rassī]{\hi{रस्सी} (rassī) — a rope}

\label{lesson:HI-C62-rope}



\begin{quote}
Before the new run of words: what did \emph{chūlhā} mean, and what did \emph{tavā} mean?
\end{quote}

\subsection*{You'll want to know: \hi{रस्सी}}
\textbf{\hi{रस्सी}} (\emph{rassī}) — \textquotedblleft{}a rope\textquotedblright{}.

From Sanskrit \hi{रश्मि} (\emph{raśmi}), which meant a cord or a rein — and also a ray of the sun. One word held both, because a ray was pictured as a thread let down from the sky.

Hindi kept the thread and let the light go: \hi{रस्सी} is plain rope, the kind that draws a bucket up out of the ground. The bright half of \hi{रश्मि} survives in names and in poetry, so the two halves of the old word now live in different registers.

\begin{grammarlens}[title={What you've built}]
The first of five things hands work with.
\end{grammarlens}

\subsection*{Your turn}
Say these aloud:

\begin{itemize}
\raggedright
  \item \emph{rassī}
  \item it once more, slowly
  \item \emph{rassī}, then \emph{sūraj}, and say what the old word held in common
\end{itemize}

\begin{tcolorbox}[breakable,colback=teal!4,colframe=teal!35!black,title={Before you move on}]
What does \hi{रस्सी} mean? (\textquotedblleft{}a rope\textquotedblright{}.) Say it once more.
\end{tcolorbox}
\section[sūī]{\hi{सूई} (sūī) — a needle}

\label{lesson:HI-C62-needle}



\begin{quote}
Before the new one: what did \emph{rassī} mean?
\end{quote}

\subsection*{You'll want to know: \hi{सूई}}
\textbf{\hi{सूई}} (\emph{sūī}) — \textquotedblleft{}a needle\textquotedblright{}.

From Sanskrit \hi{सूची} (\emph{sūcī}), 'needle', on \hi{सूच्} (\emph{sūc-}), 'to point out, to indicate'. A needle is named for the fact that it points.

The same root gives \hi{सूचना}, the word on every notice board and news bulletin for 'information' — a thing pointed out to you. A tailor's tool and a public announcement, off one small Sanskrit word about pointing.

\begin{grammarlens}[title={What you've built}]
Two.
\end{grammarlens}

\subsection*{Your turn}
Say these aloud:

\begin{itemize}
\raggedright
  \item \emph{sūī}
  \item it once more, slowly
  \item \emph{sūī}, then \emph{kapṛā}, and say which one goes through the other
\end{itemize}

\begin{tcolorbox}[breakable,colback=teal!4,colframe=teal!35!black,title={Before you move on}]
What does \hi{सूई} mean? (\textquotedblleft{}a needle\textquotedblright{}.) Say it once more.
\end{tcolorbox}
\section[ṭokrī]{\hi{टोकरी} (ṭokrī) — a basket}

\label{lesson:HI-C62-basket}



\begin{quote}
Before the new one: what did \emph{sūī} mean?
\end{quote}

\subsection*{You'll want to know: \hi{टोकरी}}
\textbf{\hi{टोकरी}} (\emph{ṭokrī}) — \textquotedblleft{}a basket\textquotedblright{}.

Here the trail stops early. \hi{टोकरी} has no agreed Sanskrit ancestor, and the retroflex \hi{ट} at the front is the sound most often credited to old contact with the Dravidian languages of the south — the same honest gap you already met under \hi{पेट}.

A \hi{टोकरी} is woven, usually from split cane or from the ribs of a leaf, and it is what \hi{सब्ज़ी} comes home from the \hi{बाज़ार} in. Light, cheap, and mended rather than replaced.

\begin{grammarlens}[title={What you've built}]
Three.
\end{grammarlens}

\subsection*{Your turn}
Say these aloud:

\begin{itemize}
\raggedright
  \item \emph{ṭokrī}
  \item it once more, slowly
  \item \emph{ṭokrī}, then \emph{sabzī} in it
\end{itemize}

\begin{tcolorbox}[breakable,colback=teal!4,colframe=teal!35!black,title={Before you move on}]
What does \hi{टोकरी} mean? (\textquotedblleft{}a basket\textquotedblright{}.) Say it once more.
\end{tcolorbox}
\section[miṭṭī]{\hi{मिट्टी} (miṭṭī) — clay, earth}

\label{lesson:HI-C62-clay}



\begin{quote}
Before the new one: what did \emph{ṭokrī} mean?
\end{quote}

\subsection*{You'll want to know: \hi{मिट्टी}}
\textbf{\hi{मिट्टी}} (\emph{miṭṭī}) — \textquotedblleft{}clay, earth\textquotedblright{}.

From Sanskrit \hi{मृत्तिका} (\emph{mṛttikā}), 'earth, clay', on \hi{मृद्} (\emph{mṛd}), 'soil'. Five syllables became two, and the meaning did not move at all.

One word does the work of two English ones. \hi{मिट्टी} is the ground under a field and the material a \hi{दीया} is pinched out of — and the \hi{दीया} you already own is the clearest case, since it is thrown away after one night and goes back to being \hi{मिट्टी}.

\begin{grammarlens}[title={What you've built}]
Four.
\end{grammarlens}

\subsection*{Your turn}
Say these aloud:

\begin{itemize}
\raggedright
  \item \emph{miṭṭī}
  \item it once more, slowly
  \item \emph{miṭṭī}, then \emph{dīyā}, and say which one the other is made of
\end{itemize}

\begin{tcolorbox}[breakable,colback=teal!4,colframe=teal!35!black,title={Before you move on}]
What does \hi{मिट्टी} mean? (\textquotedblleft{}clay, earth\textquotedblright{}.) Say it once more.
\end{tcolorbox}
\section[lakṛī]{\hi{लकड़ी} (lakṛī) — wood, a stick of firewood}

\label{lesson:HI-C62-wood}



\begin{quote}
Name the 4 words of this chapter so far.
\end{quote}

\subsection*{You'll want to know: \hi{लकड़ी}}
\textbf{\hi{लकड़ी}} (\emph{lakṛī}) — \textquotedblleft{}wood, a stick of firewood\textquotedblright{}.

Usually traced back to Sanskrit \hi{लगुड} (\emph{laguḍa}), 'a club, a staff', through a Prakrit form with a doubled consonant in the middle. A word for a stick you swing became the word for the stuff a stick is made of.

\hi{लकड़ी} is both, in fact: the material and one piece of it. A bundle of it under the arm is what feeds the \hi{चूल्हा}, which is why this word and that one turn up in the same sentence more often than any other pair in this book.

\begin{grammarlens}[title={What you've built}]
Five: \hi{रस्सी}, \hi{सूई}, \hi{टोकरी}, \hi{मिट्टी}, \hi{लकड़ी}. Enough to say what a working pair of hands is holding.
\end{grammarlens}

\subsection*{Your turn}
Say these aloud:

\begin{itemize}
\raggedright
  \item \emph{lakṛī}
  \item it once more, slowly
  \item all five in order, then \emph{merā nām} and your own name
\end{itemize}

\begin{tcolorbox}[breakable,colback=teal!4,colframe=teal!35!black,title={Before you move on}]
What does \hi{लकड़ी} mean? (\textquotedblleft{}wood, a stick of firewood\textquotedblright{}.) Say it once more.
\end{tcolorbox}
